\section{\texorpdfstring{Terminal RR Combination and the $\Delta J_n^{(0,\alpha)}$ Term}{Terminal RR Combination and the Delta J Term}}

Applying the deterministic-product decomposition of the previous subsection
separately to step sizes $\alpha$ and $2\alpha$, and writing
$B_{2\alpha}:=I-2\alpha\overline{A}$, define the chapter-local terminal
Richardson--Romberg object
\[
  \theta_n^{(\mathrm{tRR},\alpha)}
    :=2\theta_n^{(\alpha)}-\theta_n^{(2\alpha)}.
\]
It is not the PR-averaged RR estimator
$\overline{\theta}_n^{(\mathrm{RR},\alpha)}$ studied in Chapters~4--5. Its
deterministic-product decomposition has the form
\[
  \begin{aligned}
    \theta_n^{(\mathrm{tRR},\alpha)}-\theta^*
      ={}&\bigl[2B_\alpha^n-B_{2\alpha}^n\bigr](\theta_0-\theta^*) \\
       &+\bigl[2J_n^{(0,\alpha)}-J_n^{(0,2\alpha)}\bigr]
        +\bigl[2R_n^{(\alpha)}-R_n^{(2\alpha)}\bigr].
  \end{aligned}
\]
The first bracket is the deterministic transient in this convention. The
second bracket is the linearized stochastic RR difference, and the final one
is the higher-order random-product remainder. In this subsection we focus on
the zeroth-order RR difference, which we denote
\[
  \Delta J_n^{(0,\alpha)}
    :=2J_n^{(0,\alpha)}-J_n^{(0,2\alpha)},
\]
where, by definition,
\[
  J_n^{(0,\alpha)}
    =-\alpha\sum_{j=1}^n
      (I-\alpha\overline{A})^{n-j}\varepsilon(Z_j).
\]

Substituting and using the elementary identity
$X^m-Y^m=(X-Y)\sum_{i=1}^mX^{i-1}Y^{m-i}$ with $X=B_\alpha$,
$Y=B_{2\alpha}$, and $X-Y=\alpha\overline{A}$, define the deterministic
kernel
\[
  H_j^{(n)}
    :=\sum_{i=1}^{n-j}
      B_\alpha^{i-1}B_{2\alpha}^{n-j-i},
  \qquad 1\leq j\leq n-1,
\]
and set $H_n^{(n)}:=0$. Then
\[
  \begin{aligned}
    \Delta J_n^{(0,\alpha)}
      &=-2\alpha\sum_{j=1}^n
        \bigl(B_\alpha^{n-j}-B_{2\alpha}^{n-j}\bigr)\varepsilon(Z_j) \\
      &=-2\alpha^2\overline{A}\sum_{j=1}^{n-1}
        H_j^{(n)}\varepsilon(Z_j).
  \end{aligned}
\]
The left-factored matrix-power identity is legitimate here because
$I-\alpha\overline{A}$, $I-2\alpha\overline{A}$, and $\overline{A}$ are
polynomials in the same matrix $\overline{A}$ and therefore commute. The extra
factor $\alpha\overline{A}$ pulled out front is the source of the additional
$\alpha$-decay of the RR difference compared to a single LSA trajectory.
